% ============================================================
%  PlantMetBench — Technical Appendix
%  Paste this after \bibliography{...} in your main .tex file.
%  Requires: \usepackage{booktabs,xcolor,amsmath,amssymb,hyperref}
% ============================================================

\definecolor{todo}{RGB}{200,0,0}
\newcommand{\TODO}[1]{\textcolor{todo}{\textbf{[#1]}}}

% ─────────────────────────────────────────────────────────────
\appendix
% ─────────────────────────────────────────────────────────────

\section{Dataset Construction}
\label{app:dataset}

\subsection{Source Knowledge Graph}

PlantMetWiki~\TODO{cite} is a curated, RDF-based knowledge graph of plant primary
and specialised metabolism.
It encodes biochemical reactions, the metabolites they transform, the enzymes that
catalyse them, and the genes that encode those enzymes, all linked to source pathways
and to the plant organisms in which they have been characterised.
The knowledge is represented as a set of RDF triples and exported as Turtle
(\texttt{.ttl}) files.

\subsection{Graph Construction Pipeline}

We extract a heterogeneous property graph from PlantMetWiki in three stages.

\paragraph{Triple parsing.}
We parse 3.8 million RDF triples from the PlantMetWiki Turtle export using
\texttt{rdflib}.
Nodes are typed according to their RDF class (\texttt{biopax:Protein},
\texttt{biopax:BiochemicalReaction}, \texttt{biopax:SmallMolecule}, etc.),
and edges are typed by predicate.
Interaction nodes that correspond to enzymatic conversions are identified by
their \texttt{biopax:BiochemicalReaction} subclass and labelled with the
\emph{Conversion} subtype; other interaction types (transport, binding, etc.)
are retained in the graph but excluded from supervision.
Organism assignments are resolved via a supplementary NCBI taxonomy RDF file
(\texttt{taxonomy.ttl}), which maps each biological entity to its NCBI Taxon IRI.

\paragraph{Schema alignment.}
The raw RDF predicates (\texttt{participants}, \texttt{source}, \texttt{target}, \ldots)
are re-typed to a cleaner relational schema.
The central supervision edge is renamed
\texttt{(Protein, catalyzes, Interaction)}.
The reverse edge \texttt{(Interaction, catalyzed\_by, Protein)} is retained in the
raw schema but removed by default during model training (see Appendix~\ref{app:shortcuts}).

\paragraph{PyTorch Geometric conversion.}
The typed tables (\texttt{nodes.tsv}, \texttt{edges.tsv}) are loaded into a
\texttt{HeteroData} object and serialised as \texttt{heterodata.pt} using
PyTorch Geometric~\cite{Fey/Lenssen/2019}.
Pre-computed node embeddings are attached to the corresponding node stores at
this stage; node types without a domain-specific embedding receive a zero tensor
that is replaced by random features at training time to break initialisation symmetry.

\subsection{Graph Statistics}

\begin{table}[h]
\centering
\caption{Node types in the PlantMetBench graph before any shortcut removal.}
\label{tab:node_types}
\begin{tabular}{llrrl}
\toprule
Node type & Role & \# Nodes & Feature dim & Source \\
\midrule
Protein        & Enzyme candidates (query side)          & 13,682  & 960   & ESM-C (8,445 covered; rest zero) \\
Interaction    & Biochemical reactions (target side)     & 84,470  & 3,072 & MAP4 DRFP-approx \\
\quad\emph{Conversion subtype} & Enzymatic conversions (labelled) & 19,927 & 3,072 & MAP4 DRFP-approx \\
Metabolite     & Substrate / product molecules           & 17,476  & 1,024 & MAP4 \\
GeneProduct    & Encoding genes                          & 6,480   & 1,024 & PlantCaduceus \\
Pathway        & Metabolic pathways (\emph{removed})     & 2,478   & —     & — \\
Organism       & Plant species                           & 424     & 64    & random / taxonomy MDS \\
\bottomrule
\end{tabular}
\end{table}

\begin{table}[h]
\centering
\caption{Edge types in the PlantMetBench graph.
Edges marked (\emph{removed}) are absent from the message-passing graph in the
default configuration (see Appendix~\ref{app:shortcuts}).}
\label{tab:edge_types}
\begin{tabular}{llr}
\toprule
Edge type & Description & \# Edges \\
\midrule
\multicolumn{3}{l}{\textit{Supervision and catalysis}} \\
(Protein, \textbf{catalyzes}, Interaction)         & \textbf{Supervision target}                   & 11,980 \\
(Interaction, catalyzed\_by, Protein)              & Reverse supervision (\emph{removed})          & 11,980 \\
\midrule
\multicolumn{3}{l}{\textit{Organism assignment}} \\
(Protein, organism, Organism)                      & Species origin of protein                     & 13,725 \\
(GeneProduct, organism, Organism)                  & Species origin of gene                        & 6,504 \\
(Metabolite, organism, Organism)                   & Species origin of metabolite                  & 17,489 \\
\midrule
\multicolumn{3}{l}{\textit{Reaction–metabolite connectivity}} \\
(Interaction, participants, Metabolite)            & Substrate or product of reaction              & 34,793 \\
(Metabolite, is\_part\_of, Interaction)            & Reverse of above                              & 34,793 \\
(Interaction, source, Metabolite)                  & Substrate of reaction                         & 7,554 \\
(Interaction, target, Metabolite)                  & Product of reaction                           & 18,622 \\
\midrule
\multicolumn{3}{l}{\textit{Gene–protein linkage}} \\
(GeneProduct, encodes, Protein)                    & Gene encodes protein                          & 2,762 \\
(Protein, encoded\_by, GeneProduct)               & Reverse of above                              & 2,762 \\
\midrule
\multicolumn{3}{l}{\textit{Reaction hierarchy}} \\
(Interaction, is\_part\_of, Interaction)           & Sub-reaction nesting                          & 9,813 \\
(Interaction, participants, Interaction)           & Reaction–reaction participation               & 9,813 \\
\midrule
\multicolumn{3}{l}{\textit{Pathway membership (\emph{all removed})}} \\
(Interaction, is\_part\_of, Pathway)               & Reaction pathway membership                   & 29,740 \\
(Metabolite, is\_part\_of, Pathway)                & Metabolite pathway membership                 & 22,109 \\
(Protein, is\_part\_of, Pathway)                   & Protein pathway membership                    & 10,660 \\
(GeneProduct, is\_part\_of, Pathway)               & Gene pathway membership                       & 8,215 \\
(Pathway, is\_about, Pathway)                      & Cross-species pathway equivalence             & 2,478 \\
\bottomrule
\end{tabular}
\end{table}

\noindent
After removing all Pathway-related edges and \texttt{catalyzed\_by},
the message-passing graph retains 14 edge types and approximately
227,774 edges.


% ─────────────────────────────────────────────────────────────
\section{Train / Validation / Test Split}
\label{app:split}
% ─────────────────────────────────────────────────────────────

\subsection{Taxa Holdout Rationale}

We use a \emph{taxa holdout} split: entire plant species are withheld from
training and assigned exclusively to validation or test.
This design tests genuine cross-species generalisation — the model must predict
enzyme–reaction associations for organisms it has never encountered, relying on
conserved biochemistry encoded in the protein and reaction representations rather
than memorising species-specific patterns.
A random-edge split, by contrast, allows the model to look up a test reaction
in the training graph via shared pathway membership or metabolite co-occurrence,
making the task far easier than the deployment scenario.

\subsection{Organism Assignment}

Organisms are sorted by the number of positive \texttt{catalyzes} edges they
contribute.
A \emph{coverage constraint} is applied: an organism is only eligible for the
validation or test set if at least 10\% of its proteins appear in the embedding
pool (i.e.\ have a non-zero ESM-C embedding).
This prevents the evaluation from being dominated by organisms for which the
model has no feature information.

After filtering, organisms are partitioned into:
\begin{itemize}
  \item \textbf{Train:} 344 organisms contributing the majority of positive edges.
  \item \textbf{Validation:} 11 held-out organisms.
  \item \textbf{Test:} 69 held-out organisms.
\end{itemize}

\subsection{Split Statistics}

\begin{table}[h]
\centering
\caption{Positive edge counts per split.
Only \texttt{(Protein, catalyzes, Interaction)} edges where the
\texttt{Interaction} node has subtype \emph{Conversion} appear in supervision.}
\label{tab:splits}
\begin{tabular}{lrr}
\toprule
Split & \# Organisms & \# Positive edges \\
\midrule
Train      & 344 & 9,569 \\
Validation & 11  & 1,211 \\
Test       & 69  & 1,200 \\
\midrule
Total      & 424 & 11,980 \\
\bottomrule
\end{tabular}
\end{table}

\subsection{Supervision-Only Training Edges}

Within the training split, we further apply a \emph{disjoint ratio} of 0.2:
80\% of training positive edges enter the message-passing (MP) graph, while the
remaining 20\% are used only as supervision labels and are absent from the MP
graph.
This prevents the 1-hop shortcut, whereby the model could score a training pair
$(\text{Protein}_i, \text{Conversion}_j)$ by directly reading off the
\texttt{catalyzes} edge in the MP graph instead of learning a latent
compatibility score (see Appendix~\ref{app:shortcuts}).


% ─────────────────────────────────────────────────────────────
\section{Node Representations}
\label{app:embeddings}
% ─────────────────────────────────────────────────────────────

Each node type is initialised with a pre-computed, domain-specific feature vector.
Node types without a dedicated embedding receive 64-dimensional random Gaussian
features, which break initialisation symmetry and are transformed by the GNN
into task-relevant representations.

\subsection{Protein Embeddings (ESM-C)}

Protein sequences are retrieved from UniProt/RCSB.
Of the 13,682 Protein nodes in the graph, \textbf{8,445} have a resolved
amino-acid sequence and are embedded with
ESM-C~\TODO{cite}~(960-dimensional per-protein mean-pool representations).
Proteins without a sequence receive a zero vector; the model therefore sees
these proteins only through structural graph signal, and they are excluded from
the evaluation ranking pool (Appendix~\ref{app:metrics}).

\subsection{Metabolite Embeddings (MAP4)}

Metabolite SMILES strings are resolved from PubChem, MetaNetX, and KEGG for
4,529 of the 17,476 Metabolite nodes.
Each resolved metabolite is encoded with MAP4~\cite{Capecchi2020}
(MinHashed Atom-Pair fingerprint, 1,024-dimensional).
Metabolites with unresolvable SMILES (e.g.\ generalised R-group structures)
receive a zero vector.

\subsection{Gene Embeddings (PlantCaduceus)}

Gene nucleotide sequences are retrieved from NCBI Nucleotide for
1,359 of the 6,480 GeneProduct nodes.
Each resolved sequence is embedded with
PlantCaduceus~\TODO{cite}~(1,024-dimensional).

\subsection{Reaction Fingerprints (MAP4 DRFP-approximation)}
\label{app:drfp}

Biochemical reactions in the graph are represented by their substrates and products as
Metabolite nodes.
We approximate the Differential Reaction FingerPrint
(DRFP)~\cite{Probst2022} using the pre-computed MAP4 embeddings:
for a Conversion node $r$ with substrate set $S_r$ and product set $P_r$,
the reaction fingerprint is:
\begin{equation}
  \mathbf{f}_r
  = \left[
      \overline{\mathbf{m}}_{S_r}
      \;\|\;
      \overline{\mathbf{m}}_{P_r}
      \;\|\;
      \overline{\mathbf{m}}_{S_r} - \overline{\mathbf{m}}_{P_r}
    \right]
  \in \mathbb{R}^{3 \times 1024},
\end{equation}
where $\overline{\mathbf{m}}_X = \frac{1}{|X|}\sum_{x \in X}\mathbf{m}_x$
is the mean MAP4 embedding of metabolite set $X$, and $\|$ denotes
concatenation.
This yields a \textbf{3,072-dimensional} reaction fingerprint.
Of the 19,927 Conversion nodes, 13,717 have at least one resolved
metabolite embedding and thus a non-zero fingerprint; the remainder receive a
zero vector.

\subsection{Organism Features}

By default, each Organism node is initialised with a 64-dimensional random
Gaussian vector.
As an optional enhancement, Organism features can be replaced with
\emph{taxonomy-aware} features derived from NCBI taxonomic lineage:
for each of the 424 organisms, we retrieve the named-rank lineage
(phylum, class, order, family, genus) via the NCBI Taxonomy API, construct
a 702-dimensional binary multi-hot vector over all unique named-rank ancestors,
and apply classical Multidimensional Scaling (MDS / PCoA) on the pairwise
taxonomic hop-distance matrix to produce a 64-dimensional Euclidean
embedding.
All baseline results reported in the paper use the random initialisation.


% ─────────────────────────────────────────────────────────────
\section{Shortcut Analysis}
\label{app:shortcuts}
% ─────────────────────────────────────────────────────────────

A \emph{shortcut} is a structural pattern in the training graph from which a
model can infer supervision labels without learning the task-relevant biochemical
features.
We identify four shortcuts in the raw PlantMetWiki graph and remove each in the
default configuration.
Table~\ref{tab:shortcuts} summarises their effect on validation and test P-H@50.

\begin{table}[h]
\centering
\caption{Ablation of each shortcut.
Reported metric is P-H@50 at the best validation checkpoint.
\TODO{Fill in all cells from ablation runs.}}
\label{tab:shortcuts}
\begin{tabular}{lcc}
\toprule
Configuration & Val P-H@50 & Test P-H@50 \\
\midrule
Full baseline (all shortcuts removed, default)      & 11.5\%        & 6.0\% \\
\midrule
+\;Pathway co-membership (\texttt{is\_part\_of})    & \TODO{X.X\%}  & \TODO{X.X\%} \\
+\;Reverse catalysis edge (\texttt{catalyzed\_by})  & \TODO{X.X\%}  & \TODO{X.X\%} \\
Without disjoint training ratio (ratio = 0)         & \TODO{X.X\%}  & \TODO{X.X\%} \\
+\;EC one-hot features (\texttt{--ec-features})     & \TODO{X.X\%}  & \TODO{X.X\%} \\
\bottomrule
\end{tabular}
\end{table}

\subsection{Pathway Co-membership Shortcut}

PlantMetWiki encodes pathway membership via
\texttt{(Interaction, is\_part\_of, Pathway)} and
\texttt{(Protein, is\_part\_of, Pathway)} edges.
If these edges are present in the MP graph, the model can infer that a protein
$p$ and a conversion $c$ are likely to be associated simply because they co-occur
in the same pathway node — without learning any enzyme–substrate compatibility.
This is a near-trivial shortcut: two-thirds of positive pairs share at least one
pathway.

We remove all \texttt{is\_part\_of} edges and the Pathway node store (2,478 nodes,
54,800 pathway-membership edges) from the MP graph.
The pathway-membership relation is preserved in the raw \texttt{edges.tsv} file
for use in hard-negative mining and post-hoc analysis.

\subsection{Reverse Catalysis Edge}

Including \texttt{(Interaction, catalyzed\_by, Protein)} in the MP graph creates
a 2-hop shortcut: for a training pair $(p, c)$, the model can traverse
$p \xrightarrow{\texttt{catalyzes}} c \xrightarrow{\texttt{catalyzed\_by}} p$
and learn to score this neighbourhood pattern rather than the biochemical
compatibility.
We remove \texttt{catalyzed\_by} entirely from the MP graph;
the forward \texttt{catalyzes} edge is used only as a supervision label.

\subsection{1-Hop Training Supervision Shortcut}

If all positive training edges are in the MP graph, a 1-layer GNN can construct
a node embedding for $p$ that directly encodes its neighbours in $c$-space,
making training label recovery trivial.
We apply a disjoint training ratio of 0.2: 20\% of training positives are
\emph{supervision-only} (absent from the MP graph), and these are the pairs the
model is trained to score.
The remaining 80\% enter the MP graph to allow multi-hop message passing but
are never used as direct supervision labels.

\subsection{Enzyme Class (EC) Feature Shortcut}

The reaction EC number (e.g.\ EC 1.1.1.1) is a coarse functional label that
directly identifies the enzyme class responsible for a reaction.
When EC one-hot features (237-dimensional, encoding EC level-1, level-2, and
level-3 classes) are appended to Interaction node features, the model can memorise
the mapping from EC class to the set of training-organism proteins with that
annotation, achieving high validation AUC without any biochemical generalisation.
This manifests as a $\sim$80\% relative drop from validation to test P-H@50
in the taxa holdout, indicating that the model is overfitting to the distribution
of training organisms' EC annotations.
EC features are therefore \textbf{disabled by default} and provided only as an
ablation flag (\texttt{--ec-features}).


% ─────────────────────────────────────────────────────────────
\section{Evaluation Metrics}
\label{app:metrics}
% ─────────────────────────────────────────────────────────────

We use three metrics.
The primary metric directly measures the research question; the secondary metrics
provide complementary diagnostic information.

\subsection{Protein Hits at \texorpdfstring{$K$}{K} (P-H@$K$)}
\label{app:phk}

\paragraph{Definition.}
Let $\mathcal{P}_\text{emb} \subset \mathcal{P}$ be the set of proteins with
non-zero ESM-C embeddings (the \emph{ranking pool}, $|\mathcal{P}_\text{emb}| = 8{,}445$).
For each held-out positive edge $(p^*, c) \in \mathcal{E}_\text{eval}$
with $p^* \in \mathcal{P}_\text{emb}$, let
\begin{equation}
  \operatorname{rank}(p^*, c)
  = \bigl|\bigl\{q \in \mathcal{P}_\text{emb}
    : s(q, c) \geq s(p^*, c)\bigr\}\bigr|,
\end{equation}
where $s(p, c) = \mathbf{z}_p \cdot \mathbf{z}_c$ is the dot-product score
of the model embeddings.
Protein Hits at $K$ is then:
\begin{equation}
  \text{P-H@}K
  = \frac{1}{|\mathcal{E}_\text{eval}|}
    \sum_{(p^*, c) \in \mathcal{E}_\text{eval}}
    \mathbf{1}\bigl[\operatorname{rank}(p^*, c) \leq K\bigr].
\end{equation}

\paragraph{Ranking pool.}
Only proteins in $\mathcal{P}_\text{emb}$ are ranked.
Positive edges whose true catalyst $p^* \notin \mathcal{P}_\text{emb}$ are
excluded from evaluation rather than counted as misses, since the model has no
feature information about these proteins and their rank would reflect embedding
coverage, not model quality.

\paragraph{Random baseline.}
Under a uniform-random ranking over $|\mathcal{P}_\text{emb}| = 8{,}445$
proteins, the expected P-H@$K$ is $K / 8{,}445$.
For $K = 50$: random P-H@50 $\approx 0.59\%$.

\paragraph{Reported values.}
We report P-H@$K$ for $K \in \{1, 5, 10, 50\}$.
The primary metric is P-H@50; P-H@10 is a secondary ranking metric.

\subsection{Corrupt-Protein AUC and AP (CP-AUC, CP-AP)}
\label{app:cpauc}

\paragraph{Definition.}
For each positive edge $(p^*, c) \in \mathcal{E}_\text{eval}$, we sample one
\emph{corrupt protein} $\tilde{p}$ uniformly from $\mathcal{P}_\text{emb}
\setminus \{p^* \}$, ensuring $(\tilde{p}, c)$ is not a known positive.
This yields a balanced set of positive and negative pairs.
We compute:
\begin{align}
  \text{CP-AUC} &= \operatorname{AUC}\bigl(
    \{(s(p^*, c), 1)\} \cup \{(s(\tilde{p}, c), 0)\}
  \bigr), \\
  \text{CP-AP}  &= \operatorname{AP}\bigl(
    \{(s(p^*, c), 1)\} \cup \{(s(\tilde{p}, c), 0)\}
  \bigr),
\end{align}
where AUC and AP are the standard area under the ROC curve and the
average precision respectively.

\paragraph{Interpretation.}
CP-AUC and CP-AP measure whether the model scores the true enzyme above a
random protein \emph{for the same reaction}.
This is the research question in a binary, pairwise form.
The random baseline is CP-AUC $= 0.50$, CP-AP $= 0.50$.

\paragraph{Corrupt-protein pool.}
Corrupt proteins are sampled from $\mathcal{P}_\text{emb}$ (not from all proteins)
for the same reason as in P-H@$K$: proteins without embeddings are trivially
easy negatives that inflate the metric without providing information about
model quality.

\subsection{Why Not Standard Link-Prediction AUC?}

The standard link-prediction evaluation samples random $(p, c)$ pairs as
negatives, where both endpoints are chosen uniformly at random.
In the PlantMetWiki graph, Pathway co-membership provides a near-trivial
signal: two nodes drawn at random are very unlikely to share a pathway.
A model that learns only pathway membership achieves AUC $\approx 0.97$
without any biochemical generalisation, rendering this metric uninformative.
We log the random-negative AUC for historical comparability only; it
\textbf{must not be used as the primary optimisation target}.


% ─────────────────────────────────────────────────────────────
\section{Baseline Model}
\label{app:baseline}
% ─────────────────────────────────────────────────────────────

\subsection{Architecture}

The baseline model is a two-layer Heterogeneous Graph SAGE
(HeteroSAGE)~\cite{Hamilton2017} with a dot-product decoder.

\paragraph{Encoder.}
Each layer applies one \texttt{SAGEConv} per edge type (with $(-1, -1)$ lazy
input dimensions to handle differing feature sizes across node types) and
aggregates messages with a \emph{sum} aggregator across edge types
(\texttt{HeteroConv}).
Formal notation: let $\mathbf{x}_v^{(\ell)}$ be the representation of node
$v$ at layer $\ell$.
For each edge type $(u\text{-type}, r, v\text{-type})$:
\begin{equation}
  \mathbf{x}_v^{(\ell+1)}
  = \mathbf{W}^{(\ell)}_r \cdot
    \bigl[\mathbf{x}_v^{(\ell)} \;\|\; \operatorname{MEAN}_{u \in \mathcal{N}_r(v)}
    \mathbf{x}_u^{(\ell)}\bigr],
\end{equation}
and the outputs over all edge types incident to $v$ are summed.
ReLU activation and dropout ($p = 0.3$) are applied after each layer except
the last.

\paragraph{Decoder.}
The link-prediction score for an edge $(p, c)$ is:
\begin{equation}
  \hat{y}(p, c) = \mathbf{z}_p \cdot \mathbf{z}_c,
\end{equation}
where $\mathbf{z}_p, \mathbf{z}_c \in \mathbb{R}^d$ are the final-layer
embeddings of the Protein and Interaction nodes.
This simple, parameter-free decoder is effective when the GNN has already
learned discriminative embeddings.
An MLP decoder (two-layer, with hidden dimension $d$) is also provided for
ablation.

\paragraph{Initialisation.}
\texttt{SAGEConv}($-1, -1$) uses lazy parameter initialisation:
weight matrices are allocated on the first forward pass, which is triggered
by a dummy run immediately after model construction.

\subsection{Training Objective}

The model is trained with binary cross-entropy loss over sampled positive and
negative edge pairs:
\begin{equation}
  \mathcal{L}
  = -\frac{1}{|\mathcal{B}|} \sum_{(p,c,y) \in \mathcal{B}}
    \bigl[y \log \sigma(\hat{y}) + (1-y)\log(1 - \sigma(\hat{y}))\bigr],
\end{equation}
where $\sigma$ is the sigmoid function and $\mathcal{B}$ is a mini-batch of
labelled edges.

\paragraph{Negative sampling.}
For each positive training edge $(p^*, c^*)$, we sample:
\begin{itemize}
  \item $k_C = 5$ \emph{random-conversion} negatives $(p^*, c_\text{rand})$:
    same protein, random Conversion from the full 19,927-node pool.
    These train the model to distinguish the true reaction above random reactions.
  \item $k_P = 5$ \emph{random-protein} (corrupt-protein) negatives
    $(p_\text{rand}, c^*)$: random protein, same conversion.
    These train the model in the protein-ranking direction — the research question.
\end{itemize}
Both negative types exclude known positive pairs.

\subsection{Hyperparameters}

\begin{table}[h]
\centering
\caption{Baseline model hyperparameters.}
\label{tab:hparams}
\begin{tabular}{ll}
\toprule
Hyperparameter & Value \\
\midrule
\multicolumn{2}{l}{\textit{Architecture}} \\
GNN type           & HeteroSAGE (SAGEConv, aggr=sum) \\
Hidden dimension   & 128 \\
Number of layers   & 2 \\
Decoder            & Dot product \\
Dropout            & 0.3 (between layers only) \\
\midrule
\multicolumn{2}{l}{\textit{Training}} \\
Optimiser           & AdamW \\
Learning rate       & $10^{-4}$ \\
Weight decay        & $10^{-5}$ \\
Max epochs          & 200 \\
Early-stop patience & 50 (by val P-H@50) \\
Negative ratio $k_C$ & 5 \\
Negative ratio $k_P$ & 5 \\
Random seed         & 42 \\
\midrule
\multicolumn{2}{l}{\textit{Data}} \\
Disjoint train ratio  & 0.2 \\
Pathway edges         & Removed \\
Reverse edge          & Removed \\
EC features           & Off \\
Protein features      & ESM-C, 960-dim \\
Reaction features     & MAP4 DRFP-approx, 3,072-dim \\
Metabolite features   & MAP4, 1,024-dim \\
Gene features         & PlantCaduceus, 1,024-dim \\
Organism features     & Random, 64-dim \\
\bottomrule
\end{tabular}
\end{table}

\subsection{Results}

\begin{table}[h]
\centering
\caption{Baseline model results (seed = 42, best validation checkpoint).
\TODO{Add multi-seed mean $\pm$ std when 5-seed run is complete.}}
\label{tab:results}
\begin{tabular}{lccc}
\toprule
Split & P-H@50 & CP-AUC & CP-AP \\
\midrule
Random baseline  & 0.59\%  & 0.500 & 0.500 \\
\midrule
Validation       & 11.5\%  & 0.906 & 0.858 \\
Test             &  6.0\%  & 0.891 & 0.817 \\
\bottomrule
\end{tabular}
\end{table}

\noindent
The val–test gap in P-H@50 (11.5\% vs.\ 6.0\%) reflects the cross-species
generalisation challenge: the model must recognise enzyme–reaction compatibility
from protein and reaction features alone, without access to species-level context.
The CP-AUC gap is smaller (0.906 vs.\ 0.891), suggesting that pairwise protein
discrimination is more robust across species than top-$K$ ranking.
